\documentclass[12pt]{article}
\usepackage[margin=1in]{geometry}
\usepackage{tabularx,booktabs,array,siunitx}
\renewcommand{\thetable}{6.\arabic{table}}
\makeatletter\newlabel{ch:muscles}{{3}{1}}\makeatother
\begin{document}
\section*{Proposed P01: tools and demonstrations}
Researchers studying neuromechanical control need a reusable way to assemble neural controllers with human musculoskeletal and artificial-muscle models and evaluate their behavior under specified support, stimulation, and perturbation conditions. I addressed this need by developing the following models and tools from the existing simulation packages:
\begin{enumerate}
    \item A custom Li--Deng walker and controller, developed in AnimatLab and implemented in MuJoCo with SNS-Toolbox, with matched implementations of its contact-driven neural core in NumPy and Simulink.
    \item Gait2392-based neuromechanical models with alternative neural controllers, and a robot counterpart with BPA routing and the actuator-force model from Chapter~\ref{ch:muscles}.
    \item The SNS-Simscape block library and GUI for block programming neural pathways, including sensory, neuron, synapse, activation, and custom BPA-force blocks, together with coupled reflex and CPG demonstrations.
    \item A SolidWorks-derived robot assembly in Simscape Multibody coupled to a neural controller and BPA actuators using the force model from Chapter~\ref{ch:muscles}.
    \item Integration of SNS-Simscape controllers with MuJoCo through the existing MathWorks bidirectional bridge, and model-import tools for bringing musculoskeletal and CAD-derived mechanisms into the MATLAB environment.
    \item An automated curriculum and experiment workflow that combines staged parameter searches with adjustable support, neural stimuli, pathway interventions, mechanical perturbations, and recorded neural and mechanical outputs.
    \item Implementations of literature-derived circuits and selected spiking-neuron counterparts for comparing neural mechanisms and coding assumptions.
\end{enumerate}
OpenSim, MuJoCo, SNS-Toolbox, Optuna, Simulink, and the MathWorks bridge provide the underlying software; the contributions here are the model implementations, custom blocks, integrations, and evaluation procedures. Table~\ref{tab:sns_capability} presents selected functional-verification tests and proof-of-concept demonstrations illustrating the capabilities of this toolset.

\begin{table}[htbp]
\caption[Tools developed and representative demonstrations]{Tools developed in this project and representative functional-verification tests and proof-of-concept demonstrations. The examples illustrate the toolset rather than exhaustively evaluating every implementation; each result refers to the named circuit or mechanical model.}
\label{tab:sns_capability}
\small
\begin{tabularx}{\textwidth}{@{}>{\raggedright\arraybackslash}p{0.23\textwidth}*{2}{>{\raggedright\arraybackslash}X}@{}}
\toprule
\textbf{Tool developed} & \textbf{Function} & \textbf{Representative demonstration} \\
\midrule
Li--Deng/W2L model implementations & Combine the walker mechanics and neural architecture; transfer selected controller components between environments & Matched contact-driven rhythm cores reproduce \qty{1.000}{\second} periods with trace correlation at least 0.99996 \\
SNS-Simscape library and GUI & Construct neural pathways graphically and connect sensory, neural, and actuator blocks & The loaded-arm reflex demo gives \ang{2.00} sag with feedback and \ang{7.85} without it; actuator blocks agree with reference force calculations \\
Gait2392/SNS model implementations & Evaluate alternative controllers on the human reference and custom actuator/routing models & Individual-actuator tests check torque directions; selected human-model controllers produce bilateral stepping-like activity under support \\
W2L BPA-force integration & Apply the Festo force model within the coupled simulation & The 12-actuator BPA model demonstrates air stepping and retained-support standing, with peak pressure \qty{372}{\kPa} \\
SNS-Simscape/MuJoCo integration & Exchange activation commands and mechanical measurements through the MathWorks bridge & Single-muscle and control-rate tests reproduce the Python-reference sensor readback over the scored windows \\
Curriculum and experiment drivers & Automate staged tuning, support changes, neural stimuli, and mechanical perturbations & Saved stage-specific trials, candidate replays, and fixed-controller pathway and body-change comparisons demonstrate the workflow \\
Literature-circuit and spiking implementations & Test selected published neural behaviors and compare graded and spiking formulations & The reduced Rybak circuit passes 13 drive/deletion checks; selected spiking counterparts have paired implementation tests \\
\bottomrule
\end{tabularx}
\end{table}

\end{document}
